\documentclass[12pt,a4paper]{article}
\usepackage{ucebnice}
\usepackage{polynom}

\title{Základní matematické poznatky}
\author{Ondřej Škubala}
\date{\today}

\begin{document}

\maketitle
\newpage

\section{Algebraické výrazy}

S algebraickými výrazy jsme se v matematice již mnohokrát setkali.  
Jsou to například:
\[
3x + 5, \quad 2a^2b - \tfrac{1}{3}a, \quad \frac{x^2 - 1}{x+3}, \quad |x-5| + 2, \quad \tfrac{1}{3}\pi r^2 v
\]

\noindent
Každý algebraický výraz obsahuje dvě základní složky.  
\important{Konstanty} jsou čísla, která mají pevnou hodnotu a nemění se.  
\important{Proměnné} jsou písmena, která zastupují čísla z určité množiny, a tato množina se nazývá \important{obor proměnné}.  

\begin{examplebox}
Výraz $\tfrac{1}{3}\pi r^2 v$ udává objem rotačního kužele o výšce $v$ a poloměru podstavy $r$.  
\[
V = \tfrac{1}{3}\pi r^2 v
\]

\begin{itemize}
    \item Konstanty jsou zde $\tfrac{1}{3}$ a $\pi$.  
    \item Proměnné jsou $r$ a $v$, protože jejich hodnota se může měnit.  
\end{itemize}
\end{examplebox}

\noindent
S výrazy se dá dále pracovat. Často potřebujeme \important{upravit výraz}, tj. nahradit jej jiným tvarem, který má pro všechny hodnoty proměnných ve svém oboru stejnou hodnotu. Úpravami obvykle získáme jednodušší nebo vhodnější tvar, se kterým se lépe počítá.  

\begin{examplebox}
Mějme výraz pro objem kužele:
\[
V = \tfrac{1}{3}\pi r^2 v.
\]

Pokud víme, že výška kužele je $1{,}5$-násobkem poloměru, tedy $v = 1{,}5r$,  
můžeme tento vztah dosadit přímo do výrazu:  
\[
V = \tfrac{1}{3}\pi r^2 \cdot (1{,}5r) = \tfrac{1}{2}\pi r^3.
\]

Původní výraz jsme \uv{upravili} na jednodušší tvar, který už závisí pouze na poloměru $r$.  
\end{examplebox}

\noindent
Je důležité neplést si, co je konstanta a co proměnná. Například:
\begin{itemize}
    \item číslo $\pi$ je vždy konstanta,  
    \item písmeno $p$ však může být proměnnou (například značí-li délku strany čtverce).  
\end{itemize}

\noindent
Když do výrazu dosazujeme čísla, musíme vždy brát v úvahu, pro jaké hodnoty má výraz smysl. Například nesmíme dělit nulou nebo odmocňovat záporná čísla (v oboru reálných čísel). Množina všech čísel, která lze do výrazu dosadit, se nazývá \important{definiční obor výrazu}.  

\begin{examplebox}
Zjistěme, pro jaké hodnoty proměnných dávají následující výrazy smysl:

\[
a)\; \frac{x}{3-x} \qquad b)\; \frac{\sqrt{y-3}}{x^2-4}
\]

\textbf{Řešení:}

\begin{itemize}
    \item[a)] Výraz $\frac{x}{3-x}$ má ve jmenovateli $3-x$. Tento jmenovatel nesmí být nulový:  
    \[
    3 - x \neq 0 \;\;\Rightarrow\;\; x \neq 3.
    \]
    Definiční obor je tedy množina všech reálných čísel kromě $x=3$:  
    \[
    D = \mathbb{R} \setminus \{3\}.
    \]

    \item[b)] Výraz $\frac{\sqrt{y-3}}{x^2-4}$ má dvě podmínky:
    \begin{enumerate}
        \item Argument odmocniny musí být nezáporný: $y - 3 \geq 0 \;\Rightarrow\; y \geq 3$.
        \item Jmenovatel nesmí být nulový: $x^2 - 4 \neq 0 \;\Rightarrow\; x \neq \pm 2$.
    \end{enumerate}

    Definiční obor je tedy:  
    \[
    D = \{(x,y) \in \mathbb{R}^2 \;|\; y \geq 3, \; x \neq -2, \; x \neq 2 \}.
    \]
\end{itemize}
\end{examplebox}



\section{Mnohočleny (polynomy)}

Zvláštním případem algebraických výrazů jsou \important{mnohočleny}, nebo také \important{polynomy}.
Podle toho, kolik proměnných obsahují, rozlišujeme mnohočleny s jednou, či více proměnnými. 
Příkladem takových mnohočlenů jsou:
\[
2x^3y - 5xy^2 + y - 1, \quad x^3 - x^2y + 3xy^2 - y^2, \quad x^2 - 3y^2 + \sqrt{3}
\]

Jedná se o výrazy, kde proměnná vystupuje pouze s nezápornými celočíselnými exponenty.  

\begin{definitionbox}
    Obecný tvar mnohočlenu je:
    \[
    P(x) = a_0 + a_1x + a_2x^2 + \dots + a_nx^n,
    \]
    kde $a_0, a_1, \dots, a_n \in \mathbb{R}$ jsou \important{koeficienty mnohočlenu}, $a_n \neq 0$ a $n \in \mathbb{N}_0$ je \important{stupeň mnohočlenu}.  
\end{definitionbox}

\newpage
Stupeň mnohočlenu s více proměnnými udává nejvyšší součet mocnitelů proměnných v určitém členu. 
Stupeň mnohočlenu značíme $\operatorname{st}(P(x))$ nebo $\deg P(x)$. 

\begin{examplebox}
\begin{itemize}
    \item $P(x) = x^3 + 3x^2 + 3x + 1$ je mnohočlen třetího stupně.
    \item $M(x) = x^2 - \tfrac{3}{7}x + \sqrt{2}$ je mnohočlen druhého stupně.
    \item $N(x,y) = x^3 + 3x^2y + y^3$ je mnohočlen dvou proměnných, třetího stupně.
    \item $x^2 + \tfrac{1}{x} + 2$ není mnohočlen, protože obsahuje zápornou mocninu $x^{-1}$.
\end{itemize}
\end{examplebox}


Podle stupně rozlišujeme zvláštní typy mnohočlenů:
\begin{examplebox}
\begin{itemize}
    \item $A(x) = 7$ je mnohočlen 0. stupně (konstantní mnohočlen).  
    \item $B(x) = 2x + 3$ je mnohočlen 1. stupně (lineární).  
    \item $C(x) = 3x^2 + x + 6$ je mnohočlen 2. stupně (kvadratický).  
    \item $D(x) = 3x^3 - 1$ je mnohočlen 3. stupně (kubický).  
    \item $Z(x) = 0$ nazýváme \important{nulový mnohočlen}, jehož stupeň se někdy definuje jako $-\infty$.  
\end{itemize}
\end{examplebox}


\subsubsection*{Dosazování do mnohočlenů}

Do mnohočlenu můžeme dosazovat za proměnné různé číselné hodnoty.  
Pokud za proměnnou dosadíme konkrétní číslo, získáme číselný výsledek.  
Takto můžeme například ověřit, zda dva různé zápisy představují tentýž mnohočlen,  
nebo jednoduše spočítat jeho hodnotu pro zadanou proměnnou.

\begin{examplebox}
Mějme mnohočlen
\[
P(x) = 2x^3 - 5x^2 + 3x - 7.
\]

Chceme určit jeho hodnotu pro $x=2$ a $x=-1$.

\[
P(2) = 2\cdot 2^3 - 5\cdot 2^2 + 3\cdot 2 - 7 = 16 - 20 + 6 - 7 = -5,
\]

\[
P(-1) = 2\cdot (-1)^3 - 5\cdot (-1)^2 + 3\cdot (-1) - 7 = -2 - 5 - 3 - 7 = -17.
\]

Výsledky ukazují, že hodnota mnohočlenu se mění podle dosazené hodnoty proměnné.
\end{examplebox}

\newpage
\subsection{Operace s mnohočleny}
\subsubsection*{Sčítání a odčítání mnohočlenů}

Sčítání a odčítání mnohočlenů vychází z úpravy \important{stejných mocnin} proměnné.  
Koeficienty u členů se stejným exponentem se sečtou nebo odečtou, ostatní členy se jen opíší.

\begin{examplebox}
Mějme mnohočleny:
\[
A(x) = 3x^3 + 3, \qquad B(x) = x^3 + 2x^2 - 6x - 1.
\]

Součet je:
\[
A(x) + B(x) = (3x^3 + 3) + (x^3 + 2x^2 - 6x - 1) = 4x^3 + 2x^2 - 6x + 2.
\]

Rozdíl je:
\[
A(x) - B(x) = (3x^3 + 3) - (x^3 + 2x^2 - 6x - 1) = 2x^3 - 2x^2 + 6x + 4.
\]
\end{examplebox}

\subsubsection*{Opačný mnohočlen}

K danému mnohočlenu $P(x)$ definujeme \important{opačný mnohočlen} $-P(x)$ tak,
 že změníme znaménko u každého koeficientu.
Mějme:
\[
P(x) = 2x^3 - 5x^2 + 3x - 7
\]
  
Pak opačný mnohočlen je:
\[
-P(x) = -2x^3 + 5x^2 - 3x + 7.
\]

Opačný mnohočlen se využívá zejména při odčítání mnohočlenů. Například:
\[
A(x) - B(x) = A(x) + \big(-B(x)\big).
\]

Platí také, že součet mnohočlenu a jeho opačného mnohočlenu je vždy \important{nulový mnohočlen}:
\[
P(x) + (-P(x)) = 0.
\]

\begin{examplebox}
Pro $P(x) = 2x^2 - x + 4$ platí:
\[
P(x) + (-P(x)) = (2x^2 - x + 4) + (-2x^2 + x - 4) = 0.
\]
\end{examplebox}



\subsubsection*{Násobení polynomů}

Při násobení polynomů násobíme \important{každý člen jednoho mnohočlenu každým členem druhého mnohočlenu}.  
Koeficienty násobíme jako reálná čísla, exponenty u stejné proměnné sečítáme.

\begin{examplebox}
\[
A(x) = 5x^3 + 6x - 1, \qquad B(x) = x + 2
\]

\[
A(x)\cdot B(x) = (5x^3 + 6x - 1)(x+2) = 5x^4 + 10x^3 + 6x^2 + 11x - 2.
\]

Výsledkem je mnohočlen 4. stupně.
\end{examplebox}

Opakovaným násobením pak lze odvodit i umocňování mnohočlenů.
Například dvojčlen $A(x,y)=x+y$, který bychom chtěli umocnit na druhou mocninu $A^2(x,y)$ můžeme rozepsat jako součin dvou stejných závorek:  

\[
(x+y)^2 = (x+y)(x+y).
\]

Použitím distributivního zákona dostáváme:
\[
(x+y)(x+y) = x^2 + xy + yx + y^2 = x^2 + 2xy + y^2.
\]

Podobně lze odvodit i další často používané vzorce. 
Tyto vzorce nazýváme \important{algebraické vzorce} a jsou velmi užitečné při úpravách mnohočlenů.

\begin{notebox}
    \textbf{Algebraické vzorce:}
    \[
    \begin{array}{rl}
        (a+b)^2 &= a^2 + 2ab + b^2 \\[6pt]
        (a-b)^2 &= a^2 - 2ab + b^2 \\[6pt]
        (a+b)(a-b) &= a^2 - b^2 \\[6pt]
        (a+b)^3 &= a^3 + 3a^2b + 3ab^2 + b^3 \\[6pt]
        (a-b)^3 &= a^3 - 3a^2b + 3ab^2 - b^3 \\[6pt]
        a^3 + b^3 &= (a+b)\cdot (a^2 - ab + b^2) \\[6pt]
        a^3 - b^3 &= (a-b)\cdot (a^2 + ab + b^2)
    \end{array}
    \]
    Dvojčleny $(a+b)^n$ lze umocňovat užitím tzv \important{Binomické věty}, která je rozepsána v kapitole o kombinatorice.
\end{notebox}

Tyto vzorce si lze nejen pamatovat, ale také kdykoliv znovu odvodit pomocí násobení mnohočlenů.  
Právě proto je znalost násobení a distributivního zákona tak důležitá.


\subsubsection*{Dělení mnohočlenu jednočlenem}

Pokud dělíme mnohočlenu jednočlenem, vydělíme jím \important{každý člen zvlášť} a zbylé podíly sečteme.
Zároveň je dobré si uvědomit, že ne vždy musí podíl mnohočlenů vyjít opět mnohočlen.
\newpage
\begin{examplebox}
Mějme mnohočlen
\[
A(x) = 4x^6 + 8x^3 + 2x^2,
\]
který chceme vydělit jednočlenem $2x^2$:  

\[
\frac{4x^6 + 8x^3 + 2x^2}{2x^2}
= \frac{4x^6}{2x^2} + \frac{8x^3}{2x^2} + \frac{2x^2}{2x^2}
= 2x^4 + 4x + 1.
\]

Výsledkem je mnohočlen $2x^4 + 4x + 1$.

\medskip
Naopak, když vezmeme mnohočlen:
\[
A(x) = 9x^3 - 18x^2 + 27x - \tfrac{1}{9},
\]
a budeme ho dělit jednočlenem $-9x$, tak dostáváme:

\[
= \frac{9x^3}{-9x} + \frac{-18x^2}{-9x} + \frac{27x}{-9x} + \frac{-\tfrac{1}{9}}{-9x}
= -x^2 + 2x - 3 + \frac{1}{81x}.
\]

Zde již výsledkem není mnohočlen, ale racionální lomený výraz,  
protože se objevil člen $\frac{1}{81x}$, který má proměnnou ve jmenovateli a tudíž nemá mocninu v z množiny přirozených čísel, neboť $\frac{1}{81x} = (81x)^{-1}$ . 
\end{examplebox}



\subsubsection*{Dělení mnohočlenu mnohočlenem}

Pokud dělíme mnohočlenu mnohočlenem, používáme algoritmus \important{dlouhého dělení}.  
Ten postupuje podobně jako písemné dělení čísel.

\begin{enumerate}[label=\arabic*.]
    \item Vydělíme \emph{první člen dělence} prvním členem dělitele.  
    \item Výsledek zapíšeme do podílu.  
    \item Tento výsledek vynásobíme celým dělitelem a odečteme od dělence.  
    \item Pokračujeme, dokud nelze vydělit další členy.  
\end{enumerate}

Při dělení polynomu polynomem může nastat jedna ze dvou možností:
\begin{itemize}
    \item \important{Dělení beze zbytku} – pokud dělitel přesně „pasuje“ do dělence, výsledkem je opět polynom.  
    \[
    \frac{x^2 - 1}{x-1} = x+1
    \]
    \item \important{Dělení se zbytkem} – pokud dělitel není přesným dělitelem, získáme podíl a k tomu nenulový zbytek.  
    Výsledek pak zapisujeme ve tvaru
    \[
    \frac{A(x)}{B(x)} = Q(x) + \frac{R(x)}{B(x)},
    \]
    kde $A(x)$ je dělenec, $B(x)$ dělitel, $Q(x)$ podíl a $R(x)$ zbytek.  
    V příkladu výše vidíme právě tento případ.
\end{itemize}

Uvedeme si to napříkladu:

\begin{examplebox}
Určete podíl mnohočlenů
\[
(5 - 13x + 2x^3 - x^2) : (2x + 5).
\]
Oba mnohočleny musíme mít uspořádáné sestupně, čili je upravíme do tvaru:
\[
(2x^3 - x^2 - 13x + 5) : (2x + 5).
\]
Nyní první člen dělence vydělíme prvním členem dělitele, tj:

\[
2x^3 : 2x = x^2,
\]
získaným jednočlenem $x^2$ poté vynásobíme všechny člene dělitele:

\[
(2x + 5)\cdot(x^2)=2x^3 + 5x^2
\]
a vzniklý mnohočlen odečteme od dělence:
\[
(2x^3 - x^2 - 13x + 5)-(2x^3 + 5x^2) = -6x^2 - 13x + 5.
\]

Tento postup, kterým jsme dospěli k tomuto částečnému výsledku, obvykle zapisujeme následujícím způsobem:

\[
\begin{array}{rl}
    (2x^3 - \phantom{5}x^2 - 13x + 5) : (2x + 5) &= x^2 \\
    \underline{-(2x^3 + 5x^2)} \phantom{- 13x + 5) : (2x + 5) } & \\
    -6x^2 - 13x + 5 \phantom{) : (2x + 5) } & \\
\end{array}    
\]
Protože stupeň dělitele $2x + 5$ není vyšší než stupeň mnohočlenu $-6x^2 - 13x + 5$, pokračujeme v dělení,
až do momentu, kdy je mnohočlen nulový, nebo má stupeň menší, než je stupeň dělitele. Celý postup pro dělení
je z následujícího zápisu patrný:

\[
\begin{array}{rl}
    (2x^3 - \phantom{5}x^2 - 13x + 5) : (2x + 5) &= x^2 - 3x +1 \\
    \underline{-(2x^3 + 5x^2)} \phantom{- 13x + 5) : (2x + 5) } & \\
    -6x^2 - 13x \phantom{\;+ 5) : (2x + 5) } & \\
    \underline{-(-6x^2 -15x)} \phantom{ + 5) : (2x + 5) } & \\
    2x+5 \phantom{): (2x + 5) } & \\
    \underline{-(2x+5)} \phantom{)) (2x + 5) } & \\
    0 \phantom{): (2x + 5) } & \\
\end{array}    
\]
Máme tedy výsledek: pro všechna $x\in \mathbb{R}$, pro než je $2x^3 - x^2 - 13x + 5 \neq 0$, platí
$(2x^3 - \phantom{5}x^2 - 13x + 5) : (2x + 5) = x^2 - 3x +1$.
\end{examplebox}

V tomto případě, kdy dělení vyšlo bezezbytku, je podílem obou mnohočlenů mnohočlen. O správnosti se můžeme přesvědčit zkouškou,
kterou provádíme tak, že výsledný podíl vynásobíme dělitelem a musí nám vyjít dělenec.

Platí, že stupeň zbytku je vždy \important{menší než stupeň dělitele}.  

\newpage
\subsection{Rozklad mnohočlenů}

U úprav výrazů se velmi často dostaneme k situaci, kdy potřebujeme daný mnohočlen zapsat ve tvaru součinu několika jednodušších výrazů, obvykle mnohočlenů nižšího stupně.  
Této úpravě říkáme \important{rozklad mnohočlenů}.  

\noindent
Rozkládáme tak dlouho, dokud získaný tvar není dále snadno rozložitelný.  
Podobně jako u čísel hledáme prvočíselný rozklad, u mnohočlenů se snažíme najít „základní stavební kameny“ (nerozložitelné polynomy).  

\medskip
Nejčastěji se používají dvě metody:
\begin{itemize}
    \item \important{vytýkání před závorku} – z jednotlivých členů mnohočlenu vybereme jejich společný dělitel a zapíšeme jej před závorku,
    \item \important{užití algebraických vzorců} – využijeme známé vzorce, například \((a+b)^2 = a^2 + 2ab + b^2\), nebo \((a+b)(a-b) = a^2 - b^2\).
\end{itemize}

Někdy je nutné tyto metody zkombinovat: nejprve část výrazu vytkneme, a zbylý výraz pak rozložíme pomocí vhodného vzorce.

\begin{examplebox}
Rozložte mnohočlen
\[
x^5 + x^3 - x^2 - 1.
\]

\textbf{Řešení:}

Nejprve si výraz vhodně seskupíme:
\[
(x^5 + x^3) - (x^2 + 1).
\]

V první závorce můžeme vytknout $x^3$, ve druhé závorce $1$:
\[
x^3(x^2+1) - 1(x^2+1).
\]

Nyní je v obou členech společný výraz \((x^2+1)\), který vytkneme:
\[
(x^2+1)(x^3 - 1).
\]

Vidíme, že mnohočlen se rozložil na součin dvou výrazů. Druhý výraz lze ještě dále rozložit jako \emph{rozdíl krychlí}:
\[
x^3 - 1 = (x-1)(x^2+x+1).
\]

Celý rozklad tedy je:
\[
x^5 + x^3 - x^2 - 1 = (x^2+1)(x-1)(x^2+x+1).
\]
\end{examplebox}

\newpage
\begin{examplebox}
Rozložte mnohočlen
\[
7a^3bc - 14a^3b^2c^2 - 21ab^4.
\]

\textbf{Řešení:}

Podívejme se na všechny členy a hledejme jejich největší společný dělitel.  
V číselných koeficientech je společným dělitelem číslo \(7\).  
Vidíme také, že všechny členy obsahují násobek \(ab\).  

Vytkneme tedy \(7ab\):
\[
7a^3bc - 14a^3b^2c^2 - 21ab^4 = 7ab \bigl(a^2c - 2a^2bc^2 - 3b^3\bigr).
\]

V závorkách už nemají všechny členy společný činitel, a proto zde končíme.  

Celý rozklad je:
\[
7a^3bc - 14a^3b^2c^2 - 21ab^4 = 7ab(a^2c - 2a^2bc^2 - 3b^3).
\]

\medskip
Podívejme se ještě na jiný příklad:
\[
a^n + a^{n+1}.
\]

Oba členy obsahují násobek \(a^n\). Vytkneme jej:
\[
a^n + a^{n+1} = a^n(1+a).
\]

Tento výraz už dále rozložit nejde.  
\end{examplebox}


\section{Lomené výrazy}

Mezi algebraické výrazy se dále řadí také \important{lomené výrazy}.  
Jsou to výrazy ve tvaru zlomku, kde čitatel i jmenovatel jsou mnohočleny.  
Mnohočleny, které jsme dosud probírali, můžeme nazývat také \important{racionální celistvé výrazy},  
protože obsahují pouze celé mocniny proměnných. Pokud se jeden mnohočlen ocitne v čitateli  
a druhý v jmenovateli, vzniká \important{racionální lomený výraz}:

\[
\frac{A(x)}{B(x)}, \qquad B(x) \neq 0.
\]

Například:
\[
\frac{2x+3}{x-1}, \quad \frac{x^2-4}{x^2+2x+1}, \quad \frac{3a^2b}{5ab^2-1}.
\]

Stejně jako u běžných zlomků i zde platí, že \important{jmenovatel nesmí být nulový}.  
Proto je u lomených výrazů vždy nutné určit \important{definiční obor}, 
tj. množinu všech hodnot proměnné, pro které má výraz smysl, stejně jako u algebraických výrazů.

\subsubsection*{Definiční obor lomeného výrazu}

Stejně jako u běžných zlomků platí i zde, že \important{jmenovatel nesmí být nulový}.  
Proto je u každého lomeného výrazu nutné určit \important{definiční obor}, tj. množinu všech hodnot proměnné,  
pro které má daný výraz smysl.

\begin{examplebox}
Určete definiční obor výrazu
\[
\frac{x^2 - 9}{x^2 - 5x + 6}.
\]

Jmenovatel je $x^2 - 5x + 6$, což lze rozložit na $(x-2)(x-3)$.  
Výraz nemá smysl pro $x=2$ a $x=3$, protože v těchto bodech by byl jmenovatel nulový.

Definiční obor je tedy:
\[
D = \mathbb{R} \setminus \{2,3\}.
\]
\end{examplebox}

\subsubsection*{Úpravy lomených výrazů}

S lomenými výrazy pracujeme podobně jako se zlomky:  
\begin{itemize}
    \item můžeme je \important{zkracovat}, pokud mají čitatel a jmenovatel společný dělitel,
    \item můžeme je \important{sčítat a odčítat}, pokud je převedeme na společný jmenovatel,
    \item můžeme je \important{násobit a dělit} stejně jako zlomky,
    \item při úpravách musíme stále hlídat, zda se definiční obor nezmenší.
\end{itemize}

\begin{examplebox}
Zkraťte výraz
\[
\frac{x^2 - 4}{x^2 - 2x}.
\]

Čitatel lze rozložit jako $(x-2)(x+2)$, jmenovatel jako $x(x-2)$.  
Pro $x \neq 0,2$ lze výraz zkrátit:

\[
\frac{(x-2)(x+2)}{x(x-2)} = \frac{x+2}{x}, \quad x \neq 0,2.
\]
\end{examplebox}

\begin{examplebox}
Sečtěte výrazy
\[
\frac{2}{x-1} + \frac{3}{x+1}.
\]

Společný jmenovatel je $(x-1)(x+1)$. Po úpravě:

\[
\frac{2(x+1)}{(x-1)(x+1)} + \frac{3(x-1)}{(x-1)(x+1)} = \frac{2x+2 + 3x-3}{(x-1)(x+1)} = \frac{5x-1}{x^2-1}.
\]

Výraz má smysl pro všechna $x \in \mathbb{R}$ kromě $x=\pm 1$.
\end{examplebox}



\subsubsection*{Krácení a rozšiřování}

Základní vlastnost lomených výrazů je možnost \important{krácení} a \important{rozšiřování}.  
Pro výrazy $V_1, V_2, V_3$ a pro všechny hodnoty proměnných, pro které $V_2 \neq 0$ a $V_3 \neq 0$, platí:

\[
\overset{\xrightarrow{\ \large \text{krácení}\ }}%
{\underset{\xleftarrow{\ \large \text{rozšiřování}\ }}%
{\dfrac{V_1\cdot V_3}{V_2\cdot V_3} \;=\; \dfrac{V_1}{V_2}}}
\]


\begin{itemize}
    \item Přechod zleva doprava nazýváme \important{krácení} – čitatel a jmenovatel mají společný činitel $V_3$, který můžeme vykrátit.
    \item Přechod zprava doleva nazýváme \important{rozšiřování} – čitatel i jmenovatel vynásobíme stejným nenulovým činitelem $V_3$.
\end{itemize}


\subsubsection*{Sčítání a odčítání lomených výrazů}

Abychom mohli sčítat nebo odčítat dva lomené výrazy, musíme je převést na \important{společný jmenovatel}.  
Pro výrazy $V_1, V_2, V_3, V_4$ platí:
\begin{definitionbox}
    \[
    \frac{V_1}{V_2} \;+\; \frac{V_3}{V_4}
    = \frac{V_1\cdot V_4}{V_2\cdot V_4} + \frac{V_3\cdot V_2}{V_4\cdot V_2}
    = \frac{V_1\cdot V_4 + V_3\cdot V_2}{V_2\cdot V_4},
    \]

    \[
    \frac{V_1}{V_2} \;-\; \frac{V_3}{V_4}
    = \frac{V_1\cdot V_4 - V_3\cdot V_2}{V_2\cdot V_4},
    \]

    přičemž musí být $V_2 \neq 0$, $V_4 \neq 0$.
\end{definitionbox}


\begin{examplebox}
\[
\frac{2}{x-1} + \frac{3}{x+1}
= \frac{2(x+1) + 3(x-1)}{(x-1)(x+1)}
= \frac{5x-1}{x^2-1}, \quad x \neq \pm 1.
\]
\end{examplebox}


\subsubsection*{Násobení a dělení lomených výrazů}

Při násobení lomených výrazů násobíme čitatele mezi sebou a jmenovatele mezi sebou.  
Pro výrazy $V_1, V_2, V_3, V_4$ platí:

\[
\frac{V_1}{V_2} \cdot \frac{V_3}{V_4}
= \frac{V_1\cdot V_3}{V_2\cdot V_4}, 
\qquad V_2 \neq 0, \; V_4 \neq 0.
\]

Při dělení lomených výrazů dělíme tak, že první výraz násobíme převráceným druhým:  

\[
\frac{V_1}{V_2} : \frac{V_3}{V_4}
= \frac{V_1}{V_2} \cdot \frac{V_4}{V_3}
= \frac{V_1\cdot V_4}{V_2\cdot V_3},
\qquad V_2, V_3 \neq 0.
\]

\begin{examplebox}
\[
\frac{x-1}{2x} \cdot \frac{4x}{x^2-1}
= \frac{(x-1)\cdot 4x}{2x(x-1)(x+1)}
= \frac{2}{x+1}, \quad x \neq 0, \; \pm 1.
\]
\end{examplebox}


\subsubsection*{Umocňování a úpravy složených lomených výrazů}

Lomené výrazy můžeme také \important{umocňovat} a \important{zjednodušovat složené tvary}.  
Platí obdobná pravidla jako pro zlomky:

\[
\left(\frac{V_1}{V_2}\right)^n = \frac{V_1^n}{V_2^n}, \quad V_2 \neq 0, \; n \in \mathbb{N}.
\]


\begin{examplebox}
Vypočítejte
\[
\left(\frac{2a}{3b}\right)^3.
\]

Umocňujeme zvlášť čitatel i jmenovatel:
\[
\left(\frac{2a}{3b}\right)^3
= \frac{(2a)^3}{(3b)^3}
= \frac{8a^3}{27b^3}, \quad b \neq 0.
\]
\end{examplebox}


U složených lomených výrazů využíváme pravidlo pro \important{dělení zlomků}:
\[
\dfrac{\dfrac{V_1}{V_2}}{\dfrac{V_3}{V_4}}
= \frac{V_1}{V_2} : \frac{V_3}{V_4}
= \frac{V_1}{V_2} \cdot \frac{V_4}{V_3}
= \frac{V_1 \cdot V_4}{V_2 \cdot V_3},
\qquad V_2, V_3 \neq 0.
\]
\begin{examplebox}
Upravte výraz
\[
\dfrac{\dfrac{x}{y}}{\dfrac{2x}{3y}}.
\]

Použijeme pravidlo pro dělení zlomků:
\[
\dfrac{\dfrac{x}{y}}{\dfrac{2x}{3y}}
= \frac{x}{y} \cdot \frac{3y}{2x}
= \frac{3}{2}, \quad x,y \neq 0.
\]
\end{examplebox}


\newpage
\section{Vyjádření neznámé ze vzorce}

\noindent
Velmi častým úkolem v matematice, fyzice či technických oborech je \important{vyjádření neznámé ze vzorce}.  
Máme-li vzorec, který vyjadřuje vztah mezi několika veličinami, můžeme vhodnými úpravami „izolovat“ jednu z nich,  
abychom ji mohli snadno vypočítat, pokud známe ostatní.

Základním pravidlem je, že při těchto úpravách používáme pouze \important{ekvivalentní úpravy rovnic} –  
tj. úpravy, které nemění platnost rovnice. Mezi ně patří:
\begin{itemize}
    \item sčítání/odčítání stejného výrazu na obou stranách rovnice,
    \item násobení/dělení obou stran stejným nenulovým číslem či výrazem,
    \item použití inverzní operace (např. druhá odmocnina jako inverze ke druhé mocnině).
\end{itemize}

\subsubsection*{Příklad – obecný vzorec}

Mějme vzorec pro obvod obdélníku:
\[
o = 2a + 2b.
\]

Chceme vyjádřit neznámou $a$.

\[
o = 2a + 2b \quad \Rightarrow \quad o - 2b = 2a \quad \Rightarrow \quad a = \frac{o - 2b}{2}.
\]

\subsubsection*{Příklad – technický vzorec}

Z fyziky známe vztah pro hustotu:
\[
\rho = \frac{m}{V},
\]
kde $m$ je hmotnost, $V$ objem a $\rho$ hustota.

\begin{itemize}
    \item Vyjádříme-li hmotnost $m$:
    \[
    m = \rho \cdot V.
    \]

    \item Vyjádříme-li objem $V$:
    \[
    V = \frac{m}{\rho}, \quad \rho \neq 0.
    \]
\end{itemize}

\begin{examplebox}
Ohmův zákon: $U = R \cdot I$.

\begin{itemize}
    \item Proud $I$: \quad $I = \dfrac{U}{R}, \quad R \neq 0.$
    \item Odpor $R$: \quad $R = \dfrac{U}{I}, \quad I \neq 0.$
\end{itemize}
\end{examplebox}
\medskip
\begin{examplebox}
Rotační kužel (pravý kruhový kužel): objem je dán vztahem
\[
V \;=\; \frac{1}{3}\,\pi r^{2} h \quad (r>0,\; h>0).
\]
Z tohoto vztahu vyjádříme poloměr podstavy $r$ jako
\[
\\,r \;=\; \sqrt{\dfrac{3V}{\pi h}}\,\quad (V>0,\; h>0).
\]
\end{examplebox}




% -----------------------------------------------------------------
% ---------------------- P Ř Í K L A D Y --------------------------
% -----------------------------------------------------------------


\exercisepart
\setcounter{section}{1}


% =========================
% ÚLOHA 1
% =========================
\begin{exercise}[Definiční obor výrazu]
    Určete, pro které hodnoty proměnných mají následující výrazy smysl:
    \begin{tasks}(2)
        \task $\dfrac{x}{3-x}$
        \task $\dfrac{\sqrt{\,y-3\,}}{x^2 - 4}$
        \task $\dfrac{1}{x^2+y^2-1}$
        \task $\dfrac{\sqrt{\,4-z^2\,}}{z-2}$
        \task $\dfrac{\sqrt{x-1}}{\left(y-1\right)\left(y+2\right)}$
        \task $\sqrt{\dfrac{x-2}{|x-2|}} + \sqrt{x}$
    \end{tasks}
\end{exercise}


% =========================
% ÚLOHA 2
% =========================
\begin{exercise}[Přepis slovních zadání na algebraické výrazy]
    Pomocí vhodně zvolených proměnných (např. $x$, $y$) zapište výrazy, které představují:
    \begin{tasks}(1)
        \task součet trojnásobku druhé mocniny jednoho čísla a poloviny druhé odmocniny druhého čísla
        \task rozdíl druhé mocniny trojnásobku jednoho čísla a druhé odmocniny poloviny druhého čísla
        \task druhou mocninu součinu trojnásobku jednoho čísla a poloviny druhé odmocniny druhého čísla
        \task podíl absolutní hodnoty trojnásobku jednoho čísla a poloviny druhé odmocniny absolutní hodnoty druhého čísla
    \end{tasks}
\end{exercise}


% =========================
% ÚLOHA 3
% =========================
\begin{exercise}[Objem komolého jehlanu se čtvercovými podstavami]
    Přepište následující slovní popis výpočtu do algebraického výrazu pomocí vhodně zvolených proměnných:

    \smallskip
    \emph{\uv{Objem useknuté pyramidy vypočítáš tak, že změříš stranu dolního čtverce, 
    stranu horního čtverce a svislou vzdálenost obou čtverců. 
    První dvě čísla vynásob navzájem a ještě každé samo sebou, 
    tyto tři výsledky sečti a co vyjde, vynásob třetinou třetího čísla.}}

    \smallskip
    Text vychází z historického nálezu, přesněji ze staroegypského (Moscow papyrus) postupu pro výpočet objemu komolého jehlanu se čtvercovými podstavami.
\end{exercise}


\setcounter{section}{2}
% =========================
% ÚLOHA 4
% =========================
\begin{exercise}[Rozhodněte, zda se jedná o mnohočlen a určete jeho stupeň]
    Určete, zda je výraz mnohočlenem (a v jakých proměnných). Pokud ano, napište jeho stupeň.
    \begin{tasks}(2)
        \task $3x^2 + 1$
        \task $x+y$
        \task $2x^2+xy-3$
        \task $5x^2y - 2xy + 3x + 1$
        \task $xy^2 + 6xy - 3$
        \task $x^4 + 7y^3 - x^2 - 2x^{-2} + \sqrt{2y}$
        \task $x^2+3y\sqrt{x}+y^2$
        \task $x^2-\tfrac{3}{4}x+\sqrt{3}$
    \end{tasks}
\end{exercise}


% =========================
% ÚLOHA 5
% =========================
\begin{exercise}[Seřazení mnohočlenů]
    Uspořádejte členy do konvenčního pořadí (sestupně podle mocnin; pro více proměnných $x \succ y \succ z$).
    \begin{tasks}(1)
        \task $A(x) = 2x - x^2\sqrt{3} + 2 + x^4$
        \task $B(x,y) = x^2y - 2x + 4x^3 + 2xy + 2 - 3x^2$
        \task $C(x,y) = -3xy^2 + 5x^3y - 2y + x^2 - 7$
        \task $D(x,y,z) = 2z - 3x^2y + 4x^3 - y^2 + xz$
        \task $E(x) = -5 + 3x^5 - 2x^3 + x^4 - x$
        \task $F(x,y) = 7 - 2xy + y + 3x^2y^2 - x^3$
    \end{tasks}
\end{exercise}


% =========================
% ÚLOHA 6
% =========================
\begin{exercise}[Opačný polynom]
    Najděte opačný polynom k následujícím mnohočlenům:
    \begin{tasks}(1)
        \task $P_1(x) = 4x^6 - 3x^5 + 2x^4 - x^3 + 5x^2 - 7x + 9$
        \task $P_2(x,y) = 2x^3y - 5x^2y^2 + 7xy - 3x + 4y^3 - 6y^2 + 8$
        \task $P_3(x,y,z) = x^4 - 2x^3y + 3x^2y^2 - 4xy^3 + 5y^4 + 2xz - 3yz + z^2 - 1$
    \end{tasks}
\end{exercise}


% =========================
% ÚLOHA 7
% =========================
\begin{exercise}[Sčítání a odčítání mnohočlenů]
    Sečtěte a odečtěte následující mnohočleny (výsledek uveďte v konvenčním pořadí).
    \begin{tasks}(1)
        \task $A(x)=3x^2+1,\quad B(x,y)=2x^2+xy-3$
        \task $C(x,y)=x+y,\quad D(x,y)=3x^2-xy+2x-2$
        \task $E(x,y)=4x^2y - 2xy - \sqrt{3}x + 3,\quad F(x,y)=2x^2y+5xy-y^2+1$
        \task $G(x)=x^4-2x^3+x^2-1,\quad H(x)=3x^3+2x^2-4$
        \task $I(x,y,z)=2x^2y-3xz+z^2,\quad J(x,y,z)=-x^2y+4xz-2z^2+5$
    \end{tasks}
\end{exercise}


% =========================
% ÚLOHA 8
% =========================
\begin{exercise}[Zjednodušování mnohočlenů]
    Zjednodušte následující výrazy:
    \begin{tasks}(1)
        \task $2x^2 + 3x - 2 + 2(x^2 + 1) - (3x^2 - 2x + 1)$
        \task $2x^4 - 3x + x(x^3 - 2x + 2) - x^2(3x^2 - x - 2)$
        \task $(3x^2 - 2x + 1) - (2x^2 - 5x - 4) + 3(x-2)$
        \task $(x^3 - 2x^2 + x) + (2x^2 - 5x + 7) - (x^3 - x + 1)$
    \end{tasks}
\end{exercise}


% =========================
% ÚLOHA 9
% =========================
\begin{exercise}[Roznásobování mnohočlenů]
    Použijte distributivní zákon a roznásobte následující výrazy:
    \begin{tasks}(2)
        \task $(3x^2 + 4x)\cdot 5x^2$
        \task $(x^3 - 2x)(4x^2 - 7x)$
        \task $(x+y)(a+b)$
        \task $(2a+b)(3a^2+4a-b)$
        \task $c(1-c)(1-c^3)$
        \task $(6x)(x^2 - 3x + 4)$
        \task $(2x - 1)(3y + 1)(5z - 3)$
        \task $(x - x^2)(xy - x^3)(y - 5x)$
    \end{tasks}
\end{exercise}


% =========================
% ÚLOHA 10
% =========================
\begin{exercise}[Roznásobování mnohočlenů – těžší]
    Roznásobte a zjednodušte následující výrazy:
    \begin{tasks}(1)
        \task $\bigl[(7 - a)3 - 5(a - a)\bigr]\cdot 4$
        \task $5\cdot\bigl[7(x - 2y) - 6(2x - y)\bigr]$
        \task $(x^2 + y^2)\cdot x - xy\cdot(2y)$
        \task $(0{,}2y - 0{,}5b)\cdot 7a - (0{,}4a + 0{,}6b)\cdot 3b$
        \task $4(x-1)^2 (2x + 2)^3 \cdot 2 + (2x+2)^4 \cdot 2(x+1)$
        \task $(x^2 - 4)(x^2 + 4)(2x + 8) - (x^2 + 8x - 4)(4x^3)$
    \end{tasks}
\end{exercise}


% =========================
% ÚLOHA 11
% =========================
\begin{exercise}[Dělení mnohočlenu jednočlenem]
    Vydělte mnohočlen jednočlenem (zapište výsledek v konvenčním pořadí):
    \begin{tasks}(2)
        \task $(2x^2 - 4x) : 2x$
        \task $(xy^4 - x^4y^3 + y^2) : y^2$
        \task $(8a^4b^3 - 16a^3b^2 + a^2b - a) : 2a$
        \task $(-14u^2v^3 - 6u^4v^2) : -u^2v^2$
        \task $(15m^5n^2 - 10m^3n^4 + 25m^2n) : 5mn$
        \task $(12p^6q^3 - 9p^4q^2 + 3p^2q) : 3p^2q$
    \end{tasks}
\end{exercise}


% =========================
% ÚLOHA 12
% =========================
\begin{exercise}[Dělení mnohočlenu mnohočlenem]
    Rozdělte polynomy pomocí dlouhého dělení. Pokud je výsledek se zbytkem, uveďte i zbytek:
    \begin{tasks}(2)
        \task $(2x^3 + 3x^2 + x + 6) : (x + 2)$
        \task $(x^3 - 2x^2 + 1) : (x - 1)$
        \task $(-x^4 + x^3 - 4x^2 + 7x - 3) : (-x + 1)$
        \task $(x^4 + x^3 - x - 1) : (x^2 - 1)$
        \task $(9x^3 + 18x^2 - 18x - 9) : (3x - 3)$
        \task $(2x^3 - 27x^2 + 74x - 14) : (2x - 7)$
        \task $(x^4 - 5x^3 + 5x^2 - 5x - 3) : (x - 4)$
        \task $(6x^3 - 7x^2 + 5) : (2x - 1)$
    \end{tasks}
\end{exercise}

% =========================
% ÚLOHA 13
% =========================
\begin{exercise}[Dělení mnohočlenu mnohočlenem - těžší]
    Rozdělte polynomy pomocí dlouhého dělení. Pokud je výsledek se zbytkem, uveďte i zbytek:
    \begin{tasks}(1)
        \task $(14t^5 + 4t^4 - t^3 +2t^2 +3t +5):(2t^2 -1)$
        \task $(2x^7 - 5x^6 + 3x^5 + 3x^4 - 8x^3 + 6x^2 + x - 2) : (2 - x + x^3 - x^4)$
    \end{tasks}
\end{exercise}


% =========================
% ÚLOHA 14
% =========================
\begin{exercise}[Rozklad mnohočlenů]
    Rozložte následující výrazy na součin (použijte vytýkání, seskupování, vzorce apod.):
    \begin{tasks}(2)
        \task $xr - yr - x^2 + 2xy - y^2$
        \task $(a^2 + b^2 - c^2)^2 - 4a^2b^2$
        \task $9(2a - x)^2 - 4(3a - x)^2$
        \task $a^2y - aby + a^3y - ab^2y$
        % — obtížnější
        \task $(a-b)c^2 + (b-a)c^4$
        \task $(v^2 + 1)^2 - (v^2 - 2v - 1)^2$
        \task $(x+1)^4 - x^4 + 2x^2 - 1$
        \task $(x-y)^3 - x^3 + y^3$
        \task $36 - 9x^4 - 4x^2 + x^6$
        \task $(a^2 +b^2 +c^2)^2 -4a^2b^2$
    \end{tasks}
\end{exercise}


% =========================
% ÚLOHA 15
% =========================
\begin{exercise}[Rozklad mnohočlenů]
    Rozložte na součin (využijte vhodné substituce a algebraické vzorce):
    \begin{tasks}(2)
        \task $x^2 + 4x + 4$
        \task $(x + y)^4 - x^4$
        \task $x^6 - y^6$
        \task $(x + y + r)^3 - (x + y - r)^3$
        \task $(x^2 - 2x + 3)^2 - (x^2 - 2x - 3)^2$
        \task $a^4 + b^4$ \; \;\emph{ - Sofiina identita}
    \end{tasks}
\end{exercise}


\setcounter{section}{3}
% =========================
% ÚLOHA 16
% =========================
\begin{exercise}[Zjednodušování výrazů]
    Upravte výrazy a určete podmínky (proměnné, pro něž mají výrazy smysl – zejména nenulové jmenovatele):
    \begin{tasks}(1)
        \task \(\displaystyle \dfrac{a}{a-1} + \dfrac{b}{b-1} - \dfrac{a+b}{ab - a - b + 1}\)
        \task \(\displaystyle \dfrac{x^2 - 4}{x - 2}\cdot\dfrac{x+2}{x^2 - 1}\)
        \task \(\displaystyle \left(\dfrac{a}{b} - \dfrac{b}{a}\right) : \left(\dfrac{a^2 - b^2}{ab}\right)\)
        \task \(\displaystyle \left(\dfrac{x}{y} + \dfrac{y}{x}\right) : \left(\dfrac{x^2 + y^2}{xy}\right)\)
        \task \(\displaystyle \dfrac{1}{x-1} - \dfrac{1}{x+1} : \dfrac{2x}{x^2 - 1}\)
        \task \(\displaystyle \left(\dfrac{m}{m-n} + \dfrac{n}{m+n}\right) : \dfrac{2mn}{\,m^2 - n^2\,}\)
        \task $\left(\dfrac{2ax}{yz}\right) : \left( \dfrac{3bx}{ay} : \dfrac{9b^{2}z}{8a^{2}xy} \right)$
        \task $(a^2 - 1)\cdot\left(\dfrac{1}{a-1} - \dfrac{1}{a+1} - 1\right)$
        \task $\left( \dfrac{5y}{x+y} + \dfrac{5x}{\,y-x\,} + \dfrac{10xy}{\,y^{2}-x^{2}\,} \right)\cdot
               \left( \dfrac{y}{y+x} + \dfrac{x}{\,y-x\,} - \dfrac{2xy}{\,y^{2}-x^{2}\,} \right)$
        
    \end{tasks}
\end{exercise}


% =========================
% ÚLOHA 17
% =========================
\begin{exercise}[Zjednodušování výrazů]
    Upravte výrazy a určete podmínky:
    \begin{tasks}(3)
        \task \(\displaystyle \dfrac{\dfrac{a+b}{a-b}}{\dfrac{(a+b)^2}{a^2-b^2}}\)
        \task $\dfrac{1-x}{1- \frac{x}{x+1}}$
        \task $\dfrac{t- \frac{t-1}{t+1}}{1+ \frac{t\left(t-1\right)}{t+1}}$
        \task $\dfrac{a- \frac{x^2}{a}}{x- \frac{a^2}{x}}$
        \task $\dfrac{1- \frac{2y}{x} + \frac{y^2}{x^2}}{x-y}$
        \task $\dfrac{r}{1+ \frac{1}{r+1}}$
        \task $\dfrac{\frac{x+y}{x-y}+\frac{x-y}{x+y}}{\frac{1}{\left(x+y\right)^2} + \frac{1}{\left(x-y\right)^2}}$
        \task $\frac{\frac{1-x}{1-x+x^2} + \frac{1+x}{1+x+x^2}}{\frac{1+x}{1+x+x^2} - \frac{1-x}{1-x+x^2}}$
        \task $\dfrac{1}{a+ \dfrac{1}{a- \tfrac{1}{a- \frac{1}{a}}}}$
    \end{tasks}
\end{exercise}


\newpage
\setcounter{section}{4}
% =========================
% ÚLOHA 18
% =========================
\begin{exercise}[Vyjádření neznámé ze vzorce – geometrie]
    Vyjádřete uvedenou veličinu.
    \begin{tasks}(2)
        \task $c^2 = a^2 + b^2 \hfill \rightarrow a$
        \task $S = 2\pi r\,(r+v) \hfill \rightarrow r$
        \task $V = \dfrac{1}{3}\pi r^2 h \hfill \rightarrow r$
        \task $S_{\triangle} = \dfrac{(a+b)}{2}\,v \hfill \rightarrow v$
        \task $o_{\circ} = 2\pi r \hfill \rightarrow r$
        \task $V=\frac{1}{3}v\left(S_1 + \sqrt{S_1S_2} + S_2\right)\hfill \rightarrow S_2$
    \end{tasks}
\end{exercise}


% =========================
% ÚLOHA 19
% =========================
\begin{exercise}[Vyjádření neznámé ze vzorce – fyzika a chemie]
    Vyjádřete veličinu uvedenou v závorce.
    \begin{tasks}(2)
        \task $U = R\,I \hfill \rightarrow (R)$
        \task $\dfrac{1}{f} = \dfrac{1}{a} + \dfrac{1}{b} \hfill \rightarrow (b)$
        \task $c = \dfrac{n}{V} \hfill \rightarrow (n)$
        \task $E_k = \dfrac{1}{2} m v^2 \hfill \rightarrow (v)$
        \task $pV = nRT \hfill \rightarrow (T)$
        \task $F = \dfrac{G m_1 m_2}{r^2} \hfill \rightarrow (r)$
        \task $Q = m c (t_2 - t_1) \hfill \rightarrow (m)$
        \task $T = 2\pi \sqrt{\dfrac{l}{g}} \hfill \rightarrow (l)$
        \task $v = \sqrt{\dfrac{2qU}{m}} \hfill \rightarrow (m)$
        \task $R = R_0(1 + \alpha \Delta t) \hfill \rightarrow (\Delta t)$
        \task $E = h f - W \hfill \rightarrow (f)$
        \task $p = \dfrac{F}{S} + \rho g h \hfill \rightarrow (h)$
        \task $E = \sigma \varepsilon T^4 \hfill \rightarrow (T)$
        \task $\eta = 1 - \dfrac{T_2}{T_1} \hfill \rightarrow (T_2)$
        \task $\Delta G = \Delta H - T\Delta S \hfill \rightarrow (T)$
        \task $E = \dfrac{hc}{\lambda} \hfill \rightarrow (\lambda)$
        \task $A = A_0 e^{-k t} \hfill \rightarrow (k)$
        \task $N = N_0 e^{-\lambda t} \hfill \rightarrow (t)$
        \task $p = \dfrac{nRT}{V - nb} - a\!\left(\dfrac{n}{V}\right)^2 \hfill \rightarrow (V)$
        \task $\rho = \dfrac{1}{r_vT}\big[p-p_p\left(1-\frac{r_v}{r_p}\right)\big] \hfill \rightarrow (r_v)$
    \end{tasks}
\end{exercise}

\end{document}
